\documentclass[11pt]{article}
\newcommand{\ReportShortTitle}{Generated supports and amplitudes}
\newcommand{\ReportUpdated}{18 September 2026}
\usepackage[T1]{fontenc}
\usepackage[utf8]{inputenc}
\usepackage{lmodern}
\usepackage{microtype}
\usepackage[a4paper,margin=24mm,headheight=14pt]{geometry}
\usepackage{amsmath,amssymb,mathtools}
\usepackage{booktabs,tabularx,array}
\usepackage{enumitem}
\usepackage{xcolor}
\usepackage{fancyhdr}
\usepackage{titlesec}
\usepackage[hidelinks]{hyperref}

\definecolor{ink}{HTML}{18212B}
\definecolor{accent}{HTML}{135F73}
\definecolor{muted}{HTML}{5E6A73}
\definecolor{pale}{HTML}{EEF5F7}
\color{ink}
\hypersetup{colorlinks=true,linkcolor=accent,urlcolor=accent}
\setlength{\parindent}{0pt}
\setlength{\parskip}{0.55em}
\setlist{nosep,leftmargin=1.55em}
\setcounter{secnumdepth}{2}
\titleformat{\section}{\Large\bfseries\color{accent}}{\thesection}{0.6em}{}
\titleformat{\subsection}{\large\bfseries\color{ink}}{\thesubsection}{0.55em}{}
\pagestyle{fancy}
\fancyhf{}
\fancyhead[L]{\small ODMA--URA research record}
\fancyhead[R]{\small\color{muted}\ReportShortTitle}
\fancyfoot[C]{\thepage}
\renewcommand{\headrulewidth}{0.3pt}
\newcolumntype{Y}{>{\raggedright\arraybackslash}X}
\newcommand{\R}{\mathbb{R}}
\newcommand{\C}{\mathbb{C}}
\newcommand{\Z}{\mathbb{Z}}
\newcommand{\F}{\mathbb{F}}
\newcommand{\one}{\mathbf{1}}
\newcommand{\given}{\,|\,}
\newcommand{\CN}{\mathcal{CN}}
\DeclareMathOperator{\diag}{diag}
\DeclareMathOperator{\Cov}{Cov}
\DeclareMathOperator{\Var}{Var}
\DeclareMathOperator{\supp}{supp}
\DeclareMathOperator*{\argminop}{arg\,min}
\DeclareMathOperator*{\argmaxop}{arg\,max}
\newcommand{\evidence}[1]{\textcolor{accent}{\textbf{Evidence status:}} #1}
\providecommand{\ReportUpdated}{30 August 2026}
\newcommand{\ProjectReport}[4]{%
  \pagenumbering{roman}
  \begin{titlepage}
    \vspace*{18mm}
    {\large\color{accent}\textbf{ODMA--URA research record}}\\[12mm]
    {\huge\bfseries #1}\\[5mm]
    {\Large\color{muted}#2}\\[18mm]
    \begin{tabularx}{\textwidth}{@{}p{0.16\textwidth}Y@{}}
      Stage & #3\\
      Updated & \ReportUpdated\\
      Status & #4
    \end{tabularx}
    \vfill
    \textcolor{muted}{This report is a chronological research record. Measured results, interpretations, and open
    hypotheses are labelled separately. Detailed numerical ledgers remain in \texttt{results/}.}
  \end{titlepage}
  \pagenumbering{arabic}
}


\begin{document}
\ProjectReport{Generating a sparse global codebook}
{Support evidence and coordinate-wise amplitude generation}{5}
{Support tests complete. Coordinate-amplitude job 031 locally checked; HPC results pending.}

\section{What we are trying to preserve}

A $B$-bit message chooses one of $M=2^B$ columns of a common codebook:
\[
 \Phi=[\phi(w_0),\ldots,\phi(w_{M-1})]\in\R^{n\times M},\qquad
 y=\Phi a+z=\sum_{k=1}^{K}\phi(w_k)+z,\qquad \|\phi(w)\|_2^2=1.
\]
Here $a\in\Z_+^M$ counts transmitted messages and sums to $K$. Repeated messages add multiplicity.
Bits label messages, not channel amplitudes: the transmitted coordinates are real numbers in these experiments.

The useful empirical result is that a sparse global matrix can perform close to a dense matrix, whereas a particular
ODMA support restriction performs substantially worse at the same density. Job 027 at $B=12,n=256$ gave D0 mean
high-SNR PUPE $0.2935$ for dense, $0.2932$ for arbitrary support size 64, and $0.5363$ for four-mask ODMA with
the same support size. We seek inexpensive structure that preserves useful codebooks, not sparsity alone.

Storing even just $T$ nonzeros per message costs $T2^B$ values. It cannot reach $B=100$. We therefore replace stored
columns by a rule $w\mapsto\phi(w)$ and test each restriction against the explicit parent at small $B$.
D0/D1 provide a common receiver for these comparisons; neither is certified optimal. A compact forward rule does
not supply a fast inverse for a noisy sum of many messages. That separate problem remains open.

\section{Support: one selected resource in each table}

Partition $n=TR$ coordinates into $T$ tables of $R=2^r$ coordinates. For every message,
\[
 h_t(w)=A_tw\oplus b_t,\qquad A_t\in\F_2^{r\times B},\quad b_t\in\F_2^r,\qquad
 \rho_t(w)=tR+\operatorname{integer}(h_t(w)).
\]
$\F_2$ means binary arithmetic, with XOR as addition. Tables are indexed $t=0,\ldots,T-1$.
Every message uses one position per table: support size $T$, density $1/R$.
These tables are not independently chosen payload sections; each hash can depend on all $B$ bits.

One-per-table masks restrict all size-$T$ masks, and affine hashes further restrict message-to-table assignments.
Neither concession is proved free. Independent uniform table choices have expected pair overlap $T/R=T^2/n$,
matching independent uniform size-$T$ supports. Equal means do not guarantee equal tails or performance.

We impose two exact binary rank conditions:
\[
 \operatorname{rank}_{\F_2}(A_t)=r,\qquad
 \operatorname{rank}_{\F_2}\begin{bmatrix}A_0\\\vdots\\A_{T-1}\end{bmatrix}=B.
\]
The first gives equal bin populations across all messages. The second ensures distinct messages have distinct full
support tuples, requiring $Tr\ge B$. For $d=w\oplus v\ne0$, overlap in table $t$ occurs exactly when $A_td=0$.
At small $B$, enumerating differences permits selection by overlap tails. Job 031 retains the previous best-of-128
hash search. This exhaustive selection is not a large-$B$ algorithm; large-$B$ forward checks instead sample
rank-valid banks without enumeration.

\evidence{Jobs 028 and 029 returned their 36 and 20 expected runs. Under joint training at $B=14,n=256$, selected
hash minus arbitrary sparse PUPE was approximately $-0.0003/-0.0103$ for D0/D1 at $T=16$, and $+0.0026/+0.0025$
at $T=32$. The support restriction is a reasonable working choice at these tested points, not a universal theorem.}

\section{Amplitude: each coordinate has its own lookup}

Keep the support fixed. Give table $t$ a binary projection and a scalar lookup bank:
\[
 P_t\in\F_2^{J\times B},\qquad V\in\R^{T\times2^J},\qquad
 u_t(w)=\operatorname{integer}(P_tw),\qquad g_t(w)=V_{t,u_t(w)}.
\]
One message chooses $T$ scalar entries, generally from different bank columns. Gather and normalize:
\[
 g(w)\in\R^T,\qquad \alpha(w)=\frac{g(w)}{\|g(w)\|_2},\qquad
 \boxed{\phi(w)=\sum_{t=0}^{T-1}\alpha_t(w)e_{\rho_t(w)}}.
\]
Positions are distinct, so $\|\phi(w)\|_2^2=\sum_t\alpha_t(w)^2=1$ for every message.
For the exceptional $g=0$ case the implementation defines $\alpha=e_0$, also unit energy. This is not expected
with continuous initialization; if reached, actual nonzero support decreases.

For $B=14,n=256,T=32,R=8,J=4$, the support maps have shape $(32,3,14)$, amplitude maps $(32,4,14)$,
and $V$ has shape $(32,16)$: 512 real amplitude parameters rather than $32\cdot2^{14}=524{,}288$.
Storage is $O(T(r+J)B+T2^J)$; generating one sparse column costs $O(T(r+J)B+T)$ without visiting other messages.

\subsection{The equivalent explicit matrix}

For explanation only, enumerate all messages. Define one-hot lookup matrices
\[
 (S_t)_{jm}=\one\{j=h_t(w_m)\},\quad S_t\in\{0,1\}^{R\times M},\qquad
 (Q_t)_{um}=\one\{u=u_t(w_m)\},\quad Q_t\in\{0,1\}^{2^J\times M}.
\]
The raw amplitude matrix $G\in\R^{T\times M}$ has row $G_{t,:}=V_{t,:}Q_t$.
For nonzero columns let $c_m=1/\|G_{:,m}\|_2$. Then
\[
 \boxed{\Phi=
 \begin{bmatrix}S_0\diag(G_{0,:})\\ \vdots\\ S_{T-1}\diag(G_{T-1,:})\end{bmatrix}
 \diag(c)\in\R^{TR\times M}.}
\]
$S_t,Q_t,G,\Phi$ are conceptual objects, not required generator storage. The small-$B$ adapter caches row/label
indices to execute this matrix and its adjoint exactly. Current D0/D1 still score $M$ messages.

\subsection{Learning and remaining restrictions}

$A_t,b_t,P_t$ are fixed discrete choices. Fixed-amplitude runs train only the decoder; joint runs train $V$ too.
For $g\ne0$, normalization has Jacobian
\[
 \frac{\partial\alpha}{\partial g}=\frac{I_T-\alpha\alpha^\top}{\|g\|_2}.
\]
The lookup routes gradients into $V_{t,u_t(w)}$. Messages reusing an entry accumulate gradients there; XOR needs
no derivative. A stored bank column is \emph{not} generally a codeword's amplitude vector, so it must not be
projected to unit norm after updates.

Scalar reuse is still a restriction: messages sharing $P_tw$ share one raw amplitude, but their other coordinates
can differ. Each $P_t$ is a prefix of a fixed invertible $B\times B$ basis. Repeating bank entries embeds a $J$-bit
model in the $(J+1)$-bit model. At $J=B$, each table's label is bijective and arbitrary amplitudes on the chosen
supports are represented exactly. Model-class equality does not guarantee optimization reaches its best member.
At fixed $J$, changing the label maps gives a different restricted family, not a guaranteed superset of the previous one.

\section{Algebraic safeguards and initialization}

The projections satisfy the following ranks over $\F_2$:
\[
 \operatorname{rank}(P_t)=J,\qquad
 \operatorname{rank}(P_{\rm stack})=\min(TJ,B),\qquad
 \operatorname{rank}\begin{bmatrix}A_t\\P_t\end{bmatrix}=\min(B,r+J).
\]
These balance each bank's label use, remove avoidable whole-label aliases, and preserve as many amplitude labels
as possible within each hash bin. When $r+J\le B$, bin and amplitude labels are independent for uniform messages.
When $r+J>B$, full independence is impossible.

For $d=w\oplus v$, sharing both a resource and its raw amplitude entry requires $A_td=P_td=0$.
The joint kernel has dimension $B-\operatorname{rank}[A_t;P_t]$, so maximal rank minimizes this alias set.
It does not minimize every codeword correlation: different bank entries can have similar values and normalization
couples all tables.

Coordinate labels also avoid a shared-factor rank bound. With one common lookup $Q$, the raw matrix $G=VQ$
has real rank at most $\min(T,2^J)$. Different $Q_t$ need not obey that bound: $T=32,J=4$ can allow amplitude
rank 32 rather than at most 16. Nonzero-column normalization preserves this rank. The old \emph{physical} matrix
$\Phi$ need not have rank at most 16, however: table-specific support masking changes that argument.

\subsection{Why center and balance a small bank?}

Only $2^J$ Gaussian draws populate each table. At $J=4$, its sample mean has standard deviation $1/4$; its
sample energy also fluctuates. Reusing that short bank does not average away the table-specific bias. Initialize by
\[
 \mu_t=2^{-J}\sum_uV_{tu},\qquad
 s_t^2=2^{-J}\sum_u(V_{tu}-\mu_t)^2,\qquad
 V_{tu}\leftarrow(V_{tu}-\mu_t)/s_t.
\]
Each raw bank now has mean zero and mean-square one. Full-rank labels make these exact averages over uniform
messages too. Individual magnitudes still vary. Subsequent codeword normalization need not preserve exact
table-energy equality, so row energy remains a measured diagnostic.

This is an \emph{initialization}, not an ongoing constraint: learned banks can change means and energy allocation.
The motivation is finite-bank bias and observed geometry, not a claim that all optimal encoders are balanced.
A preliminary two-seed frozen-D0 probe gave PUPE about $0.4322$ for balanced coordinate $J=4$ versus $0.4306$
for unrestricted amplitudes on its restricted four-cell grid. It motivates matched training, not a non-inferiority claim.

Useful physical geometry is broader than binary ranks: balanced row energy, small harmful correlation tails,
well-conditioned active-column Gram matrices, and separation between different $K$-message sums.
Since $M\gg n$, $\Phi$ cannot have full column rank. Full row rank and isolated-message injectivity alone cannot
certify multiuser recovery.

\newpage
\section{Job 031: test amplitudes before changing the receiver}

Keep $B=14,n=256,T=32$, selected supports, global D0/D1, and paired message/noise streams fixed.
Three seeds and two decoders give 72 rows:
\begin{center}
\small
\begin{tabularx}{\textwidth}{@{}p{0.46\textwidth}Y@{}}
\toprule
Cases per seed and decoder & Question \\
\midrule
Coordinate $J=4,8,10$, balanced; fixed and joint & How much sharing preserves the unrestricted reference? \\
Coordinate $J=4$, raw; fixed and joint & Does initialization explain a low-$J$ penalty? \\
Shared-label $J=4$, raw and balanced; fixed & Are separate labels helpful beyond balancing alone? \\
Unrestricted $J=14$, raw; fixed and joint & Matched-budget parent with the same normalization/optimization. \\
\bottomrule
\end{tabularx}
\end{center}
Job 030 remains available as historical evidence, not the proposed route. All new rows use at most 120 epochs,
patience five on deterministic validation loss, and best-checkpoint restoration including epoch zero.
Spectral calibration uses a fixed starting vector. Fixed and joint runs have equal maximum data budgets but may
stop at different epochs.

Compare each family with its same-seed, same-decoder, same-training-mode parent. Report all load/SNR cells,
the 8/12-dB aggregate, paired seed gaps, stopping epochs, and pre/post geometry.
Three seeds are useful replication, not a guarantee that small gaps are absent. Do not average away high-load failures.

\evidence{Local checks cover endpoint equality, nesting, binary ranks, forward/adjoint equality, gradients,
checkpoint round trips, exact energy, and sampled $B=100$ generation. Twelve smoke paths and fixed/joint
D0/D1 mini learning runs pass. These validate execution and early learning only; HPC results are pending.}

The next decision is whether coordinate-wise scalar reuse retains the useful explicit-codebook performance.
Only afterwards should a scalable inverse be added and its separate recovery loss measured.
No sign-only, mutual-information, or new decoder construction is part of this batch.
\end{document}
